\section{Implementation}
\label{sec:implementation}

This section describes the four components that translate the architecture of
Section~\ref{sec:architecture} into running code: the SPLADE-tiny encoder,
the extraction pipeline, the executor variants, and the Strands Analyst.

% -----------------------------------------------------------------------
\subsection{SPLADE-tiny Encoder}
\label{sec:implementation:encoder}

Infon uses \texttt{naver/splade-v3-lexical}, a distilled SPLADE variant
with a 17\,MB footprint, Apache~2.0 licence, and no GPU requirement.  The
model weights are frozen; Infon does not fine-tune them.

\paragraph{Sparse activations as anchor vocabulary.}
For an input text sequence $t$, the encoder produces a sparse vector
$\mathbf{s} \in \mathbb{R}^{|\mathcal{V}|}$ following the standard SPLADE
activation formula (Section~\ref{sec:background:splade}).  Dimensions whose
activation exceeds a threshold $\tau$ become \emph{anchor vocabulary entries}.
These entries form the anchor space in which infon triples are represented and
over which the manifest pruner's bounding boxes are computed.

\paragraph{Anchor type system.}
Anchor entries are classified into four types:
\begin{itemize}
  \item \textbf{Actor} --- named entities, organisations, and products.
  \item \textbf{Relation} --- predicate labels.
  \item \textbf{Feature} --- properties and scalar attributes.
  \item \textbf{Market} --- domain-specific concept labels.
\end{itemize}
The type annotation is used by the extractor's role-assignment heuristic
(Section~\ref{sec:implementation:extraction}) and by the connective-predicate
auto-inference in the reasoner
(Section~\ref{sec:architecture:reasoner}).

% -----------------------------------------------------------------------
\subsection{Extraction Pipeline}
\label{sec:implementation:extraction}

\paragraph{Role assignment.}
Given a sentence and its SPLADE anchor vector, the extractor assigns each
activated anchor to a subject, predicate, or object role.  Role assignment
combines the anchor-type prior (actors favour subject/object positions;
relations favour predicate positions) with a lightweight dependency heuristic
that parses head--modifier relationships in the token sequence.

\paragraph{Dual-partition actor reranking.}
A known failure mode in simple actor extraction is direction inversion: the
extractor may emit \texttt{nvidia/partner/b200} (treating a product as the
object partner) instead of \texttt{nvidia/partner/tsmc} (the correct
actor-to-actor link).  This arises because actors and non-actor objects are
drawn from the same activation pool, and activation magnitude does not encode
syntactic direction.

The dual-partition fix separates the candidate pool into two subsets:
\emph{actors as subjects} and \emph{actors as objects}.  A word-order penalty
is then applied to each subset independently, preferring anchors whose surface
position in the sentence is consistent with the role being filled.  After this
fix, the extraction pipeline correctly produces \texttt{nvidia/partner/tsmc},
and the actor-to-actor real-data evaluation accuracy increases from 80\% to
100\%.

\paragraph{Coverage diagnostics.}
\texttt{extraction\_report()} analyses the cassettes in a store and flags:
missing anchors (schema entries that appear in no extracted infon), overfit
objects (anchors that appear as objects far more frequently than as subjects
or predicates, suggesting role confusion), and sparse cassettes (documents
from which fewer infons were extracted than expected given document length).
The report is computed within 1\,second of ingest and is automatically
returned by the \texttt{Analyst}'s \texttt{ingest} tool.

% -----------------------------------------------------------------------
\subsection{Executor Variants}
\label{sec:implementation:executors}

\paragraph{SyncExecutor.}
A single-threaded executor that processes documents sequentially.  Intended
for testing, small corpora, and interactive use.  Measured: ingesting 48
documents with \texttt{SyncExecutor} completes in approximately 0.5\,seconds.

\paragraph{ProcessExecutor.}
A multi-process executor that fans out ingest across a pool of workers, each
running an independent SPLADE-tiny instance.  Intended for production
workloads of 100 or more documents.  Workers are stateless; the cassette
format's idempotency guarantee means that any worker crash can be retried
without coordination.

\paragraph{Lambda-compatible container.}
The same \texttt{InfonStore} API operates in an AWS Lambda container.
\texttt{fsspec} routes S3 URIs transparently, so the only code difference
between a local run and a Lambda deployment is the store path:
\texttt{InfonStore(``s3://bucket/prefix'')} instead of
\texttt{InfonStore(``./local/path'')}.  Because the PyTorch dependency
exceeds the Lambda layer limit of 250\,MB, a container image is provided
(\texttt{lambda\_container.py}) that packages the full dependency set into an
ECR-deployable image.

% -----------------------------------------------------------------------
\subsection{Strands Analyst}
\label{sec:implementation:analyst}

The \texttt{Analyst} is a Strands-powered agent that wraps an
\texttt{InfonStore} and exposes it to conversational queries.  It provides
9 tools:

\begin{itemize}
  \item \texttt{set\_schema(ontology\_json)} --- activate an ontology
    definition, passed as a dict; no temporary files.
  \item \texttt{ingest(documents\_json)} --- delta ingest, followed by an
    automatic \texttt{extraction\_report} that surfaces coverage issues.
  \item \texttt{reingest()} --- re-extract all documents under the currently
    active schema.
  \item \texttt{extraction\_report()} --- return coverage diagnostics, dead
    anchors, and overfit objects for the current store.
  \item \texttt{ask(s, p, o, polarity)} --- single-claim verdict with cited
    source infons.
  \item \texttt{connect(source, target, max\_hops)} --- multi-hop MCTS chain
    verdict (\textsc{supports} / \textsc{refutes} / \textsc{nei}).
  \item \texttt{any\_of(source, targets\_json)} --- one tree walk, $N$ targets,
    ranked verdicts.
  \item \texttt{record\_finding(title, body, tags)} --- persist a finding to
    \texttt{<root>/findings/} for cross-session memory.
  \item \texttt{list\_findings(tag, limit)} --- retrieve prior findings by
    tag.
\end{itemize}

\paragraph{System prompt constraints.}
The agent's system prompt enforces four non-negotiables: (1)~never output a
verdict without citing the \texttt{ask()} sources that produced it;
(2)~when $m_\Theta > 0.7$, explicitly report that the corpus does not answer
the question; (3)~call \texttt{list\_findings} at the start of each session
so that prior investigations are visible; (4)~run \texttt{extraction\_report}
before answering on any newly ingested store.

\paragraph{Schema bootstrap loop.}
On a cold-start corpus with no prior schema, the \texttt{Analyst} can propose
an initial ontology, ingest the documents, inspect the extraction report, and
refine the schema iteratively.  On a 15-document AI-chip corpus starting from
no schema, this loop reaches 100\% document coverage in 2 iterations.

\paragraph{One-minute start.}
The following code snippet illustrates the full API: ingest, single-claim
query, multi-hop connectivity, one-of-many reachability, and conversational
use.

\begin{lstlisting}[language=Python]
from infon import InfonStore, Query, Analyst

store = InfonStore("./data/chips", schema_path="schemas/auto.json")

store.ingest(documents)               # delta append; idempotent
report = store.extraction_report()    # coverage diagnostics
print(report.summary())               # flags missing anchors, overfit objects

v = store.ask(Query().where(subject="toyota",
                             predicate="invest",
                             object="solid_state"))
print(v.label, v.mass.supports, v.mass.theta)  # SUPPORTS 0.53 0.29

# Multi-hop:
v = store.connect("toyota", "catl")   # chain -> SUPPORTS via panasonic
vs = store.any_of(
    "toyota",
    {"catl", "lg", "samsung"},        # one tree walk, N verdicts
)

# Conversational:
a = Analyst(store)
print(a("Does Toyota partner with Panasonic?"))  # agent translates, cites
\end{lstlisting}

For S3 deployment, only the store path changes:

\begin{lstlisting}[language=Python]
store = InfonStore("s3://acme/chips", schema_path="schemas/auto.json")
# same API, fsspec handles the rest
\end{lstlisting}
